\documentclass[a4paper,11pt]{article}

% ---------------------------------------------------------------
% Sprache und Typographie
% ---------------------------------------------------------------
\usepackage{fontspec}
\usepackage[ngerman,provide=*]{babel}
\usepackage{csquotes}

\usepackage[a4paper,margin=2.5cm,headheight=14pt]{geometry}
\usepackage{setspace}
\onehalfspacing
\usepackage{parskip}
\usepackage{booktabs}
\usepackage{tabularx}
\usepackage{array}
\usepackage{amsmath}
\usepackage{enumitem}
\setlist[itemize]{leftmargin=*,topsep=2pt,itemsep=2pt}
\setlist[enumerate]{leftmargin=*,topsep=2pt,itemsep=2pt}

\usepackage{fancyhdr}
\pagestyle{fancy}
\fancyhf{}
\renewcommand{\headrulewidth}{0.3pt}
\renewcommand{\footrulewidth}{0pt}
\fancyhead[L]{\footnotesize MTA -- Sitzung 3: Vom Verfahren zur Interpretation}
\fancyhead[R]{\footnotesize Seite \thepage}
\fancyfoot[C]{\footnotesize Prof.\ Dr.\ Christian Papilloud · Institut für Soziologie · MLU Halle-Wittenberg}

\usepackage{titlesec}
\titleformat{\section}{\large\bfseries}{\thesection.}{0.6em}{}
\titleformat{\subsection}{\normalsize\bfseries}{\thesubsection}{0.6em}{}
\titleformat{\subsubsection}{\normalsize\itshape}{}{0pt}{}
\titlespacing*{\section}{0pt}{1.5em}{0.5em}
\titlespacing*{\subsection}{0pt}{1.0em}{0.3em}

\usepackage{xcolor}
\definecolor{boxgrau}{gray}{0.95}
\usepackage[most]{tcolorbox}
\newtcolorbox{merkbox}[1][]{
  colback=boxgrau,colframe=black!40,boxrule=0.4pt,
  arc=2pt,left=8pt,right=8pt,top=4pt,bottom=4pt,
  fonttitle=\bfseries,#1
}

\usepackage{hyperref}
\hypersetup{
  colorlinks=true,
  linkcolor=black!70,
  urlcolor=black!60,
  citecolor=black!60,
  pdftitle={MTA Sitzung 3 -- Handout},
  pdfauthor={Christian Papilloud},
}

\title{\textbf{Vom Verfahren zur Interpretation}\\[4pt]
\large Sitzung 3 -- Vorverarbeitung, Metadaten\\und ein kritischer Blick auf Topic-Modelle\\[8pt]
\normalsize Begleitendes Handout}
\author{Prof.\ Dr.\ Christian Papilloud\\
\small Institut für Soziologie -- Martin-Luther-Universität Halle-Wittenberg}
\date{Master Soziologie -- 2026}

\begin{document}

\maketitle
\thispagestyle{fancy}

\vspace{-1em}

% ===============================================================
\section{Worum es heute geht}

In den ersten beiden Sitzungen haben wir Topic-Modelle (TM)
\emph{methodologisch} positioniert und \emph{technisch} verstanden:
quantitativ basiert, qualitativ einsetzbar, näher an Mayring als an
Oevermann; vom Grundprinzip der Matrixfaktorisierung bis zu LDA, NMF,
BERTopic und LLMs.

Die heutige Sitzung schließt den theoretischen Teil ab. Sie behandelt
das, was zwischen \emph{Forschungsfrage} und \emph{Algorithmus} liegt --
und das, was \emph{nach} dem Algorithmus kommt. Vier Blöcke:

\begin{enumerate}
\item \textbf{Vor dem Algorithmus.} Was ist ein Dokument? Wie reinigt
man eine Textquelle? Wie funktioniert der TF-IDF-Score?
\item \textbf{Metadaten.} Wie binden wir Variablen wie Zeit, Autor:in
oder Quelle in die Analyse ein? Drei Strategien.
\item \textbf{Nach dem Algorithmus.} Der kritische Blick: zwei
verbreitete Erwartungen, die Topic-Modelle \emph{nicht} erfüllen
können.
\item \textbf{Forschungspraxis.} Was sich für unsere Arbeit ändert.
Übergang zur eigentlichen Praxis mit MTA.
\end{enumerate}

\begin{merkbox}
\textbf{Eine wichtige Botschaft vorab.} Die wichtigste Operation eines
Topic-Modell-Verfahrens findet nicht im Algorithmus statt, sondern
\emph{davor} (in der Vorbereitung der Daten) und \emph{danach} (in der
Interpretation). Der Algorithmus selbst ist relativ einfach. Was ein
gutes von einem schlechten Topic-Modell unterscheidet, sind die
Entscheidungen, die der/die Forschende um den Algorithmus herum trifft.
\end{merkbox}

% ===============================================================
\section{Was \emph{vor} dem Algorithmus geschieht}

\subsection{Was ist ein \emph{Dokument}?}

Diese Frage klingt trivial. Sie ist es nicht. Bevor wir ein
Topic-Modell rechnen können, müssen wir entscheiden, was \emph{eine
Beobachtungseinheit} ist.

Beispiele:
\begin{itemize}
\item Bei einer Sammlung von Zeitungsartikeln: jeder Artikel = ein
Dokument? Oder jeder \emph{Absatz} = ein Dokument? Die Wahl ändert
das Ergebnis fundamental.
\item Bei einem qualitativen Interview: das ganze Transkript = ein
Dokument? Oder jede Antwort einzeln? Oder eine \emph{Sequenz}
zusammengehöriger Antworten?
\item Bei einer Sammlung von Werken: jedes Buch = ein Dokument? Oder
jedes Kapitel?
\end{itemize}

Die Wahl der \emph{Dokumentgröße} ist eine \textbf{methodologische
Entscheidung}, keine technische. Sie hängt von der Forschungsfrage ab.
Zu \emph{kleine} Dokumente führen dazu, dass die Maschine zu wenig
Kontext für jedes \enquote{Dokument} sieht. Zu \emph{große} Dokumente
führen dazu, dass mehrere Themen \emph{in einem} Dokument vermischt
werden und vom Modell nicht mehr unterschieden werden können.

\subsection{Stoppwörter und Bereinigung}

Topic-Modelle arbeiten auf \emph{Wortbasis}. Die häufigsten Wörter
einer Sprache (\emph{der}, \emph{die}, \emph{das}, \emph{und},
\emph{aber}) tragen aber keine inhaltliche Information. Sie würden,
\emph{wenn man sie behält}, in allen topics oben erscheinen und die
Ergebnisse unbrauchbar machen.

\paragraph{Stoppwörter} sind diese semantisch leeren Wörter. Sie
werden vor der Analyse entfernt. Listen sprachspezifischer Stoppwörter
sind frei verfügbar (für Deutsch z.\,B.\ in MTA enthalten). Aber:
\textbf{Sie können diese Liste an Ihr Korpus anpassen}. Wenn alle Ihre
Dokumente das Wort \emph{Universität} enthalten, weil sie alle aus
einem Universitäts-Newsletter stammen, dann ist \emph{Universität}
funktional ein Stoppwort und gehört auf die Liste.

\paragraph{Weitere Reinigungsschritte:}
\begin{itemize}
\item Zeichensetzung, Zahlen, Sonderzeichen entfernen.
\item In transkribierten Interviews: \enquote{äh}, \enquote{hmmm},
\enquote{[lacht]} usw.\ entfernen.
\item Sehr seltene Wörter entfernen (z.\,B.\ Wörter, die in weniger
als 2--3 Dokumenten vorkommen). Sie bringen Rauschen, keine
Information.
\end{itemize}

\subsection{TF-IDF -- den \enquote{Wortwert} messen}

Wenn das Vokabular zu groß ist, ist es sinnvoll, nur die
\emph{aussagekräftigsten} Wörter zu behalten. Aber wie misst man
\emph{Aussagekraft}? Eine etablierte Heuristik: der
\textbf{TF-IDF-Score}.

\paragraph{Idee.} Ein Wort ist aussagekräftig, wenn:
\begin{itemize}
\item es \emph{häufig genug} im Korpus vorkommt (sonst sehen wir es
zu selten);
\item aber \emph{nicht in jedem} Dokument vorkommt (sonst trennt es
keine Dokumente voneinander).
\end{itemize}

Formal: für jedes Wort $i$ berechnet man die \emph{Term Frequency}
$\mathrm{tf}_i$ (wie oft das Wort vorkommt) und die \emph{Inverse
Document Frequency} $\mathrm{idf}_i = \log_{10}(N / n_i)$, wobei $N$
die Gesamtzahl der Dokumente ist und $n_i$ die Anzahl der Dokumente,
in denen das Wort vorkommt.

\begin{merkbox}
\textbf{Beispiel aus dem Buch.} Im Reuters-Korpus (806\,791 Texte)
haben die Wörter \emph{Versuch} und \emph{Versicherung} eine ähnliche
Korpushäufigkeit (10\,422 vs.\ 10\,440). Aber \emph{Versuch} steht in
8\,760 Dokumenten, \emph{Versicherung} nur in 3\,997. Der IDF-Wert
von \emph{Versicherung} ist daher höher: das Wort ist auf wenige
Dokumente \emph{konzentriert} -- ein Hinweis auf Aussagekraft.
\end{merkbox}

\subsection{Was MTA \emph{nicht} tut: Lemmatisierung, NER, Phrasen}

In der computer-linguistischen Forschung gibt es weitere Operationen
der Vorverarbeitung, die für TM hilfreich sein \emph{können}:

\paragraph{Lemmatisierung.} Wörter werden auf ihre Grundform
zurückgeführt: \emph{ging, gingen, gegangen} $\to$ \emph{gehen};
\emph{Häuser, Hauses} $\to$ \emph{Haus}. Reduziert das Vokabular und
verbessert oft die topics. \emph{Stemming} (ältere Variante) ist
weniger empfehlenswert: es kürzt Wörter mechanisch und kann Bedeutungen
verschmelzen.

\paragraph{Named Entity Recognition (NER).} Erkennt Eigennamen
(\emph{Berlin, Merkel, FAZ}) und kann sie als zusammengesetzte
\enquote{Wörter} behandeln. Erhöht die inhaltliche Dichte der topics.

\paragraph{Phrasen.} Wortgruppen, die oft zusammen vorkommen
(\emph{soziale Frage}, \emph{Martin-Luther-Universität}), werden als
\emph{ein} Token behandelt. Inhaltlich oft präziser als die
Einzelwörter.

\textbf{In MTA werden diese Schritte nicht standardmäßig durchgeführt.}
Das ist eine bewusste Entscheidung: für die meisten qualitativen
Analysen mit moderat großen Korpora sind sie nicht entscheidend, und
sie würden die Installation des Programms deutlich komplizierter
machen. Wenn Ihr Projekt sie verlangt, müssen Sie Ihre Texte
\emph{vor} der MTA-Analyse mit einem externen Tool vorbereiten
(z.\,B.\ spaCy oder TreeTagger für Lemmatisierung).

% ===============================================================
\section{Metadaten}

Bisher haben wir Texte als \emph{reine Wortansammlungen} behandelt.
Aber unsere Dokumente haben in der Regel \emph{Metadaten}: Datum,
Autor:in, Quelle, Zielgruppe, Genre usw.

Es gibt \textbf{drei Strategien}, um diese Metadaten in eine
Topic-Modell-Analyse einzubinden.

\subsection{Strategie 1 -- Spaltung}

\paragraph{Idee.} Das Korpus wird nach Metadatenwerten in Teilkorpora
zerlegt. Man rechnet ein separates Topic-Modell für jede Teilmenge.

\paragraph{Beispiel.} Sie haben 2000 Artikel: 1000 aus der FAZ, 1000
aus der taz. Sie rechnen zwei Topic-Modelle, eines pro Zeitung, und
vergleichen sie.

\paragraph{Vorteil.} Die topics können sich pro Gruppe \emph{wirklich}
unterscheiden -- jede Gruppe bekommt ihre eigene Wortverteilung.

\paragraph{Nachteil.} Die Topic-Listen sind nicht direkt vergleichbar.
\enquote{Topic 3 bei FAZ} und \enquote{Topic 3 bei taz} sind zwei
verschiedene Dinge, die manuell in Beziehung gesetzt werden müssen.

\paragraph{In MTA.} Möglich: führen Sie MTA zweimal aus, einmal pro
Teilkorpus.

\subsection{Strategie 2 -- Hinzufügung}

\paragraph{Idee.} Die Metadaten werden \emph{nach} der Topic-Inferenz
in die Analyse eingebunden. Man rechnet \emph{ein} Topic-Modell auf
dem gesamten Korpus und betrachtet dann, wie sich die
Topic-Verteilungen \emph{nach Gruppen} unterscheiden.

\paragraph{Beispiel.} Sie rechnen 8 topics auf allen Migrations-Artikeln.
Dann fragen Sie: kommt Topic 3 (\enquote{Wirtschaftliche Aspekte}) bei
der FAZ häufiger vor als bei der taz?

\paragraph{Vorteil.} Flexibel und statistisch gut kontrollierbar. Man
kann statistische Tests anwenden -- in MTA z.\,B.\ den Welch-t-Test
mit Benjamini-Hochberg-Korrektur (im Menü \enquote{Group Comparison}).

\paragraph{Nachteil.} Die topics selbst sind \emph{nicht} durch die
Metadaten beeinflusst. Wenn die Metadaten die topics
\emph{strukturieren} sollten, geht das verloren.

\subsection{Strategie 3 -- Lernergänzung}

\paragraph{Idee.} Das Topic-Modell selbst wird so erweitert, dass die
Metadaten \emph{Teil der statistischen Modellierung} werden. Beispiele
sind \emph{Dynamic Topic Models} (Blei \& Lafferty 2006) für
Zeitserien oder \emph{Author Topic Models} (Rosen-Zvi et al.\ 2004)
für Autor:innen-Effekte.

\paragraph{Vorteil.} Die strukturelle Wirkung der Metadaten wird
direkt im Modell abgebildet.

\paragraph{Nachteil.} Komplex zu implementieren, eingeschränkte
Software-Verfügbarkeit, schwerer zu interpretieren.

\paragraph{In MTA nicht verfügbar.} Die Empfehlung von Papilloud \&
Hinneburg (2018, S.~35): Strategie 2 (Hinzufügung) ist meistens
vorzuziehen, wenn die Metadaten nicht zwingend strukturell modelliert
werden müssen.

\subsection{Synopse der drei Strategien}

\smallskip

\footnotesize
\begin{tabularx}{\textwidth}{@{}p{0.20\textwidth}XXX@{}}
\toprule
& \textbf{Spaltung} & \textbf{Hinzufügung} & \textbf{Lernergänzung} \\
\midrule
Topic-Inhalt
  & kann pro Gruppe variieren
  & gemeinsam für alle Gruppen
  & strukturell beeinflusst \\
\addlinespace[3pt]
Interpretation
  & gruppenweise
  & vergleichend
  & vereinheitlicht \\
\addlinespace[3pt]
Statistische Tests
  & schwer möglich
  & \textbf{ja} (Welch, Mann-Whitney, BH)
  & im Modell integriert \\
\addlinespace[3pt]
Verfügbar in MTA
  & ja (per zwei Durchläufe)
  & ja (Group Comparison)
  & nein \\
\bottomrule
\end{tabularx}
\normalsize

% ===============================================================
\section{Der kritische Blick: zwei falsche Erwartungen}

Wir kommen zum Kapitel 8 des Buches (\emph{Rück- und Ausblick}). Zwei
verbreitete Erwartungen werden dort behandelt. Beide werden Sie in der
Literatur und in der Praxis immer wieder hören. Die erste muss man
\emph{nuancieren} – es ist nicht einfach falsch, von Brücken zwischen
qualitativer und quantitativer Forschung zu sprechen, wenn man Topic-Modelle
einsetzt; es ist allerdings falsch, daraus eine Auflösung der Unterscheidung
abzuleiten. Die zweite muss man \emph{klar zurückweisen}.

\subsection{Erwartung 1 -- \enquote{Topic-Modelle als Brücke zwischen
Quali und Quanti?}}

Das Argument lautet meistens: Da Topic-Modelle Statistik und Mathematik
einsetzen, \emph{lösen sie} das alte Paar \emph{qualitativ vs.\
quantitativ} \emph{vollständig auf}. Sie wären eine Art \emph{Brücke}
zwischen beiden Lagern.

\paragraph{Die starke Version ist falsch} (Buch, S.~152) -- die
schwächere Version ist begründet. Wir behandeln beide Aspekte
nacheinander.

\paragraph{Was an der starken Version nicht stimmt.} Topic-Modelle
setzen die Arbeit an \emph{Texten} voraus -- also am \emph{Sinn} von
Begriffen und Sprachen. Diese Daten lassen sich nicht auf eine
einzige Bedeutung reduzieren. Sie können nicht rein quantitativ
interpretiert werden, \enquote{wie eine Variable}. Topic-Modelle
reduzieren die Komplexität von Textquellen \emph{nicht}; sie sind
\emph{im Gegenteil} dazu gemacht, die \emph{Multiperspektivität} einer
Textquelle nicht zu reduzieren, sondern zu \emph{strukturieren}. Statt
die Komplexität eines Textes zu reduzieren, laden TM dazu ein, sie zu
\emph{erkennen} und zu \emph{vertiefen}. Die hermeneutische Lektüre
bleibt unverzichtbar.

\paragraph{Was an der schwächeren Version richtig ist.} Wenn man die
Behauptung der \emph{Auflösung} der Unterscheidung nicht teilt, bleibt
ein methodisch ernst zu nehmender Punkt: Topic-Modelle eröffnen
\emph{Anschlüsse} zwischen qualitativer und quantitativer Forschung,
und zwar in einer Weise, die ältere qualitative Verfahren (Oevermann,
Mayring) so nicht erlauben. Drei konkrete Beispiele, die diese These
stützen:

\begin{itemize}
\item \textbf{Formale Parallele zur Faktorenanalyse.} Die Struktur
eines Topic-Modells ähnelt jener einer \emph{Faktorenanalyse}: In
beiden Fällen wird eine hochdimensionale Datenmatrix (hier: die
Dokument-Wort-Matrix) auf wenige \emph{latente Dimensionen}
reduziert. Topics entsprechen formal Faktoren, Dokument-Topic-Gewichte
entsprechen Faktorenwerten. Diese Parallele ist nicht nur
oberflächlich -- NMF beispielsweise ist ein klassisches
Verfahren der \emph{linearen Algebra}, das schon vor seiner
Anwendung auf Texte in der quantitativen Datenanalyse verwendet
wurde.

\item \textbf{Mittelwerttests auf Dokument-Topic-Gewichten.} In MTA
selbst führen Sie auf der Seite \enquote{Group comparison}
\emph{Welch-t-Tests} mit Benjamini-Hochberg-Korrektur durch, um zu
prüfen, ob sich Dokumentgruppen (z.\,B. nach Geschlecht, Altersgruppe,
Quelle) signifikant in ihren Topic-Gewichten unterscheiden. Das ist
ein klassisch \emph{quantitatives Verfahren}, angewandt auf Daten
\emph{qualitativer Herkunft} (Texte). Diese Operation ist nur
möglich, weil Topic-Modelle ihre Ergebnisse als numerische Gewichte
ausgeben, die ein Mittelwert-Vergleich überhaupt erst zulässt.

\item \textbf{Topics als Inputs für eine Hauptkomponentenanalyse.}
Die Topic-Gewichte selbst können in einem weiteren Schritt einer
\emph{Hauptkomponentenanalyse} (PCA) unterzogen werden, um Korpora
in einem stark reduzierten Raum zu vergleichen oder Cluster von
Dokumenten zu identifizieren. Diese Verkettung
(Texte~$\to$~Topics~$\to$~Komponenten) ist ein Beispiel für eine
Forschungspraxis, die qualitative und quantitative Operationen
\emph{ineinander verschränkt}, ohne sie zu verwechseln.
\end{itemize}

\begin{merkbox}
\textbf{Pointe.} Topic-Modelle sind \emph{kein} Brückenschlag, der
die Unterscheidung qualitativ/quantitativ \emph{aufhebt}. Aber sie
sind ein Verfahren, das \emph{Anschlüsse} zwischen beiden Traditionen
ermöglicht -- und zwar auf eine Art, die ältere qualitative Verfahren
so nicht erlauben. Die qualitative Arbeit am \emph{Sinn} bleibt
unverzichtbar, ebenso wie die methodische Reflexion über das, was die
quantitative Operation des Algorithmus eigentlich tut.
\end{merkbox}

\subsection{Erwartung 2 -- \enquote{Topic-Modelle reduzieren die
Interpretation auf ein Minimum}}

Das Argument: Da der Algorithmus die Kategorien selbst findet, müssen
wir nur noch die Ergebnisse \enquote{lesen}. Im Vergleich zu Oevermann
oder Mayring soll TM eine Befreiung von der mühsamen Codierarbeit sein.

\paragraph{Warum das nicht stimmt} (Buch, S.~152--153):

Die Ergebnisse eines Topic-Modells hängen \emph{stark} ab von:
\begin{itemize}
\item der Forschungsfrage (was wir suchen);
\item den theoretischen Rahmenannahmen;
\item den Vorverarbeitungsentscheidungen (Stoppwörter,
Dokumentgröße, Lemmatisierung);
\item der Wahl der Anzahl topics $K$;
\item der Wahl von LDA oder NMF (und in der neuen Welt: BERTopic
oder LLM);
\item dem Korpus selbst und seinen Auswahlentscheidungen.
\end{itemize}

\paragraph{Konsequenz.} Die topics sind \emph{Hypothesen}, keine
Befunde. Ihre Bedeutung muss durch \emph{Lektüre} der besten
Dokumente, durch \emph{Vergleich} mit Theorie und durch
\emph{Diskussion} erarbeitet werden. Die menschliche
Interpretationsarbeit verschwindet nicht; sie \emph{verlagert sich}.

\subsection{Pflichtlektüre: Schmidt 2012}

Eine wichtige Referenz, die das Buch zitiert: \textbf{Benjamin
Schmidt} (2012), der anhand historischer und geographischer Daten
zeigt, wie Topic-Modelle zu \emph{falschen Ergebnissen} führen können,
wenn sie \enquote{blind} angewendet werden.

Schmidts Forderung: Topic-Modelle in den Geistes- und
Sozialwissenschaften müssen \textbf{mit Vorsicht und kritischer
Reflexivität} eingesetzt werden. Konkret heißt das:

\begin{enumerate}
\item Topics nicht für bare Münze nehmen -- immer mit den
\emph{besten Dokumenten} konfrontieren (welche Dokumente haben in
diesem Topic das höchste Gewicht?).
\item Mehrere Modelle laufen lassen (verschiedene $K$, LDA \emph{und}
NMF) und Konsistenz prüfen.
\item Ergebnisse mit \emph{unabhängigen} Quellen oder Kennzahlen
vergleichen.
\item Die \emph{theoretische Rahmung} der Analyse explizit machen.
\end{enumerate}

% ===============================================================
\section{Wo bleibt der/die Forschende?}

Am Ende der ersten Sitzung haben wir eine Frage offen gelassen:

\begin{merkbox}
Eine Sozialwissenschaftlerin, die mit MTA arbeitet, codiert nicht
selbst -- sie liest die Codierung einer Maschine. Verschiebt das die
Subjektivität der Interpretation, oder verlagert es sie nur an einen
anderen Ort?
\end{merkbox}

Die Antwort lautet jetzt: \emph{ja, es verlagert sie}. Aber die
Verlagerung hat zwei interessante Eigenschaften.

\paragraph{Sie wird sichtbarer.} Alle Vorverarbeitungs- und
Modellierungsentscheidungen sind dokumentierbar: welche Stoppwörter,
welche Dokumentgröße, welche $K$-Werte getestet wurden, welches Modell
gewählt wurde, warum. Bei einer klassischen qualitativen Inhaltsanalyse
sind diese Entscheidungen oft impliziter.

\paragraph{Sie wird prüfbar.} Eine andere Forscherin kann die Analyse
mit denselben Parametern \emph{wiederholen} und prüfen, ob sie zu den
gleichen topics kommt. Das ist eine Form von Reproduzierbarkeit, die
qualitative Methoden traditionell schwer erreichen.

Diese beiden Eigenschaften sind \emph{nicht trivial}. Aber sie ersetzen
\emph{nicht} die hermeneutische Arbeit am Material. Sie verändern nur
die Bedingungen, unter denen diese Arbeit stattfindet.

% ===============================================================
\section{Was sich an der Forschungspraxis ändert}

Topic-Modelle verändern den \emph{Umgang} mit Texten in der
sozialwissenschaftlichen Forschung. Aus dem Buch (S.~153):

\begin{itemize}
\item Wir können Texte \emph{nicht mehr nur sequentiell, sondern auch
parallel} lesen. Statt Buch nach Buch oder Interview nach Interview
durchzugehen, können wir Hunderte oder Tausende von Dokumenten
gleichzeitig betrachten.

\item Wir können \emph{sehr große} Sammlungen berücksichtigen --
ganze Werke einer Autorin, alle Interviews eines Projekts, alle
Artikel zu einem Thema über zehn Jahre. Das verändert das Kriterium
\emph{dessen, was relevante Literatur zu einem Thema ist}.

\item Wir bekommen einen \emph{neuen Typ von Hypothesenproduktion}:
Hypothesen, die wir ohne den maschinellen Blick auf die ganze
Sammlung nicht hätten formulieren können.
\end{itemize}

Diese Veränderungen haben einen \emph{Preis}: die Verlagerung der
Subjektivität, die Notwendigkeit zur Dokumentation aller Schritte,
die Verpflichtung zur kritischen Reflexion. Das Buch nennt das die
\enquote{Veränderung der Forschungspraxis} (S.~153) -- und sie ist
nicht reversibel.

% ===============================================================
\section{Zusammenfassung der drei Sitzungen}

Drei Sitzungen, drei Antworten auf drei Fragen:

\begin{enumerate}
\item \textbf{Wovon sind Topic-Modelle eine Methode?} Sie sind eine
\emph{qualitative} Methode der Sozialwissenschaften -- aber näher an
Mayrings qualitativer Inhaltsanalyse als an Oevermanns objektiver
Hermeneutik. Sie teilen die drei Grundannahmen (Nicht-Zufälligkeit
des Sinns, Strukturiertheit, Rolle der Deutung), unterscheiden sich
aber im Begriff von \enquote{Algorithmus} und in der Rolle der
Kontextvariation.

\item \textbf{Wie funktionieren sie?} Durch Faktorisierung der
Dokument-Wort-Matrix in zwei kleinere Matrizen. Die zwei klassischen
Verfahren sind \emph{LDA} (probabilistisch) und \emph{NMF}
(algebraisch). Sie gehören zur Familie der Mixed-Membership-Modelle.
Neuere Verfahren (BERTopic, LLMs) ergänzen das Bild -- aber sie
ersetzen die klassischen Verfahren nicht, weder methodologisch noch
praktisch.

\item \textbf{Was können sie nicht?} Sie heben die Unterscheidung
qualitativ/quantitativ nicht auf -- aber sie eröffnen
\emph{Anschlüsse} zwischen beiden Traditionen (formale Nähe zur
Faktorenanalyse, Mittelwerttests auf Dokument-Topic-Gewichten,
mögliche Weiterverwendung in einer PCA). Sie reduzieren nicht die
Interpretationsarbeit; sie sind keine neutralen Werkzeuge. Die
Hauptarbeit liegt \emph{vor} dem Algorithmus (Vorverarbeitung) und
\emph{nach} ihm (Interpretation, kritische Reflexion).
\end{enumerate}

\begin{merkbox}
\textbf{Zum Mitnehmen für die Praxis.} Wenn Sie in den nächsten Wochen
mit MTA arbeiten, denken Sie an die drei Sitzungen zurück. Sie werden
sehen, dass die Software Sie immer wieder zwingt, Entscheidungen zu
treffen, die theoretisch motiviert sein müssen. Das ist kein Mangel
der Software -- das ist die Pointe der Methode.
\end{merkbox}

% ===============================================================
\section*{Zitierte Literatur}
\addcontentsline{toc}{section}{Zitierte Literatur}

\footnotesize

\textsc{Grundtext} der Sitzung:

Papilloud, C., \& Hinneburg, A.\ (2018). \emph{Einführung in die
qualitative Analyse von Texten mit Topic-Modellen}. Wiesbaden:
Springer. \emph{Bes.\ Kap.~2.5 (Vorverarbeitung), Kap.~2.5.5 (Metadaten),
Kap.~8 (Rück- und Ausblick).}

\bigskip

\textsc{Vorverarbeitung und TF-IDF:}

Boyd-Graber, J., Mimno, D., \& Newman, D.\ (2014). Care and Feeding
of Topic Models. In: E.~M.\ Airoldi et al.\ (Hrsg.), \emph{Handbook
of Mixed Membership Models and Their Applications}. CRC Press.

Manning, C.\ D., Raghavan, P., \& Schütze, H.\ (2008). \emph{Introduction
to Information Retrieval}. Cambridge UP.

Salton, G., Wong, A., \& Yang, C. S.\ (1975). A vector space model
for automatic indexing. \emph{Communications of the ACM} 18(11):
613--620.

Schofield, A., \& Mimno, D.\ (2016). Comparing apples to apple: The
effects of stemmers on topic models. \emph{TACL} 4: 287--300.

\bigskip

\textsc{Metadaten-Modelle:}

Blei, D.\ M., \& Lafferty, J.\ D.\ (2006). Dynamic Topic Models.
In: \emph{ICML 2006}, 113--120.

Mimno, D., \& McCallum, A.\ (2008). Topic models conditioned on
arbitrary features with Dirichlet-multinomial regression.
\emph{UAI 2008}.

Roberts, M.\ E., Stewart, B.\ M., \& Tingley, D.\ (2014). stm:
R package for Structural Topic Models. \emph{Journal of Statistical
Software}.

Rosen-Zvi, M., et al.\ (2004). The author-topic model for authors
and documents. \emph{UAI 2004}.

\bigskip

\textsc{Kritische Reflexion:}

Mimno, D., \& Blei, D.\ (2011). Bayesian Checking for Topic Models.
\emph{EMNLP 2011}.

Schmidt, B.\ M.\ (2012). Words alone: Dismantling topic models in
the humanities. \emph{Journal of Digital Humanities} 2(1).

Gelman, A., et al.\ (2013). \emph{Bayesian Data Analysis}.
Chapman \& Hall/CRC.

\bigskip

\textsc{Zur Vertiefung:}

Mohr, J.\ W., \& Bogdanov, P.\ (2013). Introduction -- Topic models:
What they are and why they matter. \emph{Poetics} 41(6): 545--569.

DiMaggio, P., Nag, M., \& Blei, D.\ (2013). Exploiting affinities
between topic modeling and the sociological perspective on culture.
\emph{Poetics} 41(6): 570--606.

\end{document}
